%%%%%%%%%%%%%%%%%%%%%%%%%%%%%%%%%%%%%%%%%%%%%%%%%%%%%%%%%%%%%%%%%%%%%%%%%%%%%%%%
%% 01_Estudo_Viabilidade_Requisitos.tex
%% Estudo de Viabilidade e Engenharia de Requisitos da Infraestrutura
%%%%%%%%%%%%%%%%%%%%%%%%%%%%%%%%%%%%%%%%%%%%%%%%%%%%%%%%%%%%%%%%%%%%%%%%%%%%%%%%

\chapter{Estudo de Viabilidade e Requisitos}
\label{cap:viabilidade_requisitos}

\section{Contexto Institucional e Objeto do Projeto}

Este projeto apresenta a especificação técnica para implantação da infraestrutura de redes e do serviço de geoprocessamento voltado ao grupo de pesquisa em Modelagem Hidrológica em Bacias Hidrográficas Nordestinas, vinculado ao Centro de Ciências Exatas e Tecnologia (CCET) da Universidade Federal de Sergipe (UFS), sob coordenação do Prof. Dr. Alcigeimes Batista Celeste.

O grupo de pesquisa desenvolve estudos sobre os impactos causados por resíduos sólidos e práticas antrópicas nos sistemas de drenagem urbana e esgotamento sanitário, além de atuar no planejamento e gerenciamento de recursos hídricos. As atividades demandam armazenamento relacional, análise espacial contínua e publicação cartográfica web. Para atender a essas necessidades, adotou-se a plataforma NextGIS Web integrada ao banco de dados relacional e espacial PostgreSQL com extensão PostGIS.

A infraestrutura física é dimensionada para operação local nas dependências do campus universitário, englobando cabeamento estruturado, racks, switches gerenciáveis e servidores dedicados. De forma complementar, uma instância em nuvem pública (AWS EC2) é empregada para o treinamento operacional dos pesquisadores e homologação do serviço.

\section{Delimitação do Espaço Físico e Escopo de Instalação}

A instalação física contempla dois prédios do campus universitário, com base nas diretrizes do Edital de Seleção Pública de Fornecedores nº 011/2025 da Fundação Espírito-Santense de Tecnologia (FEST):
\begin{enumerate}
    \item \textbf{Prédio CT-09 (Sala 101)}: Área destinada às bancadas de trabalho dos pesquisadores e Sala de Servidores, abrigando o rack principal de telecomunicações;
    \item \textbf{Prédio CT-12 (Sala 201)}: Laboratório de Geoprocessamento, abrigando o rack secundário de distribuição local.
\end{enumerate}

\section{Levantamento Quantitativo de Densidade de Pontos}

A distribuição física de pontos de rede atende às estações de trabalho, aos pontos de acesso sem fio (\textit{Access Points}) e aos enlaces de telecomunicações. A Tabela~\ref{tab:densidade_pontos} apresenta o quantitativo dimensionado para os dois ambientes.

\begin{table}[htbp]
\centering
\small
\caption{Quantitativo de pontos de telecomunicações e certificação física.}
\label{tab:densidade_pontos}
\begin{tabularx}{\textwidth}{lYYr}
\toprule
\textbf{Edifício / Sala} & \textbf{Finalidade} & \textbf{Especificação do Ponto} & \textbf{Qtd} \\
\midrule
CT-09 (Sala 101) & Estações dos Pesquisadores & Módulo Keystone RJ-45 Cat6 para mobiliário & 49 \\
CT-09 (Sala 101) & Rede Sem Fio (Ponto 901) & Tomada RJ-45 Cat6 em caixa de teto & 1 \\
CT-09 (Sala 101) & Certificação de Enlace & Teste avulso com certificador de canal & 47 \\
\midrule
CT-12 (Sala 201) & Laboratório de Geoprocessamento & Módulo Keystone RJ-45 Cat6 para mobiliário & 48 \\
CT-12 (Sala 201) & Rede Sem Fio (Ponto 1201) & Tomada RJ-45 Cat6 em caixa de teto & 1 \\
CT-12 (Sala 201) & Certificação de Enlace & Teste avulso com certificador de canal & 49 \\
\bottomrule
\end{tabularx}
\end{table}

O projeto adota cabeamento de par trançado Categoria 6 (Cat6) com suporte a frequências de até 250 MHz para todos os pontos horizontais. A interligação entre os edifícios utiliza cabo de fibra óptica monomodo com transceivers Gigabit Ethernet.

\section{Engenharia de Requisitos}

Os requisitos do sistema foram estruturados de forma atômica e verificável, com foco prioritário na conectividade, comutação, roteamento e segurança da infraestrutura de rede, integrando o suporte à plataforma de software.

\subsection{Requisitos Funcionais de Rede e Infraestrutura}

A Tabela~\ref{tab:requisitos_funcionais} detalha os requisitos funcionais estabelecidos para os ativos de comunicação, segmentação de tráfego e serviços de suporte.

\begin{table}[htbp]
\centering
\small
\caption{Especificação atômica dos requisitos funcionais do projeto.}
\label{tab:requisitos_funcionais}
\begin{tabularx}{\textwidth}{|l|Y|l|l|}
\hline
\rowcolor{tableheaderbeige}
\textbf{ID} & \textbf{Declaração Atômica do Requisito} & \textbf{Verificação} & \textbf{Prioridade} \\ \hline
RF-01 & A rede local deve segregar o tráfego em quatro VLANs independentes (Gerência, Servidores, Pesquisa e Wi-Fi). & Teste & Obrigatório \\ \hline
RF-02 & O roteador interno deve encaminhar pacotes inter-VLAN por meio de subinterfaces tronco 802.1Q. & Teste & Obrigatório \\ \hline
RF-03 & O servidor DHCP dedicado (\texttt{SRV-DHCP}) na VLAN 20 deve fornecer concessões dinâmicas de IP para as redes de Pesquisa e Wi-Fi, com o roteador interno operando como agente de repasse (\textit{DHCP Relay}). & Teste & Obrigatório \\ \hline
RF-04 & O firewall físico corporativo (Cisco ASA 5506-X) deve realizar inspeção de estado de pacotes (\textit{stateful inspection}) entre as zonas perimetrais \textit{outside} e \textit{inside}. & Teste & Obrigatório \\ \hline
RF-05 & O roteador de borda (\texttt{router\_border}) deve executar tradução de endereços por sobrecarga (\textit{NAT Overload}) para acesso à rede do provedor institucional (\texttt{ISP-UFS}). & Inspeção & Obrigatório \\ \hline
RF-06 & O roteador interno deve aplicar Listas de Controle de Acesso (ACL) para bloquear tráfego não autorizado entre sub-redes. & Teste & Obrigatório \\ \hline
RF-07 & Os pontos de acesso sem fio devem suportar o padrão Wi-Fi 6 com alimentação elétrica via portas PoE+ (IEEE 802.3at). & Inspeção & Essencial \\ \hline
RF-08 & A comunicação entre os prédios CT-09 e CT-12 deve utilizar enlace de fibra óptica Gigabit Ethernet (1000BASE-LX). & Teste & Obrigatório \\ \hline
RF-09 & A infraestrutura deve hospedar o serviço NextGIS Web acessível via HTTP/HTTPS na porta 80 e 443 TCP. & Demonstração & Obrigatório \\ \hline
RF-10 & A infraestrutura deve fornecer instância do banco PostgreSQL/PostGIS para persistência de dados espaciais. & Teste & Obrigatório \\ \hline
\end{tabularx}
\end{table}

\subsection{Requisitos Não-Funcionais de Rede e Desempenho}

A Tabela~\ref{tab:requisitos_nao_funcionais} especifica os limites físicos, métricas de capacidade de transmissão e restrições de segurança da rede.

\begin{table}[htbp]
\centering
\small
\caption{Especificação atômica dos requisitos não-funcionais do projeto.}
\label{tab:requisitos_nao_funcionais}
\begin{tabularx}{\textwidth}{|l|Y|l|l|}
\hline
\rowcolor{tableheaderbeige}
\textbf{ID} & \textbf{Declaração Atômica do Requisito} & \textbf{Critério de Aceitação} & \textbf{Prioridade} \\ \hline
RNF-01 & As portas dos switches da camada de acesso devem operar na velocidade de 1000 Mbps (Gigabit Ethernet). & Teste de porta & Obrigatório \\ \hline
RNF-02 & O comprimento do enlace permanente de cabeamento horizontal não deve ultrapassar 90 metros. & Medição com certificador & Obrigatório \\ \hline
RNF-03 & O cabeamento de par trançado deve operar com largura de banda nominal de até 250 MHz (Categoria 6). & Teste de certificação & Obrigatório \\ \hline
RNF-04 & A rede deve bloquear tentativas de conexão direta de estações de pesquisa à porta TCP 5432 do PostgreSQL. & Auditoria de ACL & Obrigatório \\ \hline
RNF-05 & O acesso administrativo por SSH e HTTPS aos ativos de rede deve ser restrito ao IP do administrador. & Teste de conexão & Obrigatório \\ \hline
RNF-06 & O sistema de energia ininterrupta (nobreak) deve manter os equipamentos do rack operantes por 30 minutos. & Teste de autonomia & Essencial \\ \hline
RNF-07 & O enlace óptico de campus deve apresentar atenuação total de inserção inferior a 3 dB no canal monomodo. & Medição com OTDR & Essencial \\ \hline
\end{tabularx}
\end{table}
